\documentclass{situationist-preprint}

\title{Using Persons in Ordinary Transactions}
\msid{ETH-21-b}
\msfamily{ETH-21}
\msdiscipline{Normative and Applied Ethics}
\mstype{Consequence}
\mskeywords{formula of humanity; mere means; exchange; bargaining; exploitation; deception}

\begin{document}
\maketitle

\begin{abstract}
ETH-21-a argues that the formula of humanity forbids only treating persons merely as means, and analyzes the qualifier in terms of reasonable rejection. We derive what this implies for ordinary transactions: purchases, employment, bargaining and the use of services provided by people one never meets. Consensual exchange without deception is permissible use; hard bargaining is permissible unless its terms exploit a party's lack of acceptable alternatives; the use of persons whose agency one does not engage at all is permissible when it runs through institutions they could not reasonably reject; and deception about the terms of a transaction always treats the deceived merely as a means, whatever its effect on her welfare.
\end{abstract}

\section{Framework}

We use the definitions of ETH-21-a. An agent \emph{uses} a person if her maxim's end is to be brought about in part by means of his body, agency, beliefs or abilities (Definition~2.1); he \emph{could reasonably reject} a way of being used if he has a ground for objecting that does not depend on his pursuing a wrongful end and is not outweighed by others' grounds for objecting to the alternatives (Definition~2.2); and he is treated \emph{merely as a means} if he is used in a way he could reasonably reject (Definition~2.3). Proposition~3.1 of ETH-21-a states that in ordinary transactions to which the parties consent each uses the other and treats the other as an end. That proposition covers the central case; what follows extends it to four kinds of case in which ordinary transactions depart from it.

\section{Consequences}

\begin{proposition}[Honest exchange]\label{p:exchange}
In an exchange whose terms are known to both parties and accepted by each without coercion, neither treats the other merely as a means, whatever the motives of either.
\end{proposition}

\begin{proof}
Each uses the other, since each acts on a maxim whose end is achieved by the other's agency. A party who knows and freely accepts the terms has no ground to object to being used on those terms that is not outweighed by the other's ground for objecting to the alternative of no exchange, unless the terms are exploitative, a case treated in Proposition~\ref{p:bargain}. So the use is not reasonably rejectable, and by Definition~2.3 of ETH-21-a neither is treated merely as a means. The parties' motives do not enter Definition~2.3: a shopkeeper indifferent to his customers' good treats them as ends if he deals with them on terms they could not reasonably reject.
\end{proof}

The last point departs from readings on which treating someone as an end requires an attitude of concern for him. On the account of ETH-21-a, the formula constrains the ways persons are used, and an indifferent but honest trader satisfies it. \citet{ONeill1985} reaches a similar verdict by asking whether the other party could consent to the transaction, which on our account is approximately the condition that he could not reasonably reject it.

\begin{proposition}[Bargaining]\label{p:bargain}
Driving a hard bargain does not treat the other party merely as a means unless the terms take advantage of his lack of acceptable alternatives to secure a share of the gains he could reasonably reject.
\end{proposition}

\begin{proof}
A party who could walk away to an acceptable alternative and accepts a hard bargain has no reasonable ground for rejecting it, since others could reasonably reject terms that bound them to give way in every negotiation. A party without acceptable alternatives has a ground for objecting to terms that exploit his position, namely that the other's gain comes from his vulnerability rather than from the value of what he contributes, and this ground does not depend on his pursuing a wrongful end. Where it is not outweighed, the terms treat him merely as a means.
\end{proof}

Proposition~\ref{p:bargain} locates the wrong of exploitation in the formula of humanity without supposing that every unequal exchange is wrong. It agrees with the view that exploitation can be mutually advantageous and consensual and still wrong \citep{Wertheimer1996}, and it explains why: consent to terms one has no acceptable alternative to accepting does not remove a reasonable ground for rejecting them.

\begin{proposition}[Remote use]\label{p:remote}
Using the work of persons one never deals with, as when one buys goods made by workers elsewhere, treats them merely as means if and only if the arrangements through which their work reaches one use them in ways they could reasonably reject.
\end{proposition}

\begin{proof}
The buyer's maxim achieves its end in part by means of the workers' abilities, so she uses them. Whether she uses them merely as means is fixed by Definition~2.3, which concerns the way they are used, and the way is fixed by the arrangements through which their work reaches her. If those arrangements involve forced labor or deception of the workers, the workers could reasonably reject them; if not, not.
\end{proof}

Proposition~\ref{p:remote} implies that the formula can be violated by ordinary purchases, but only through the arrangements, and that the buyer's responsibility depends, as for other wrongs, on what she could reasonably know about them.

\begin{proposition}[Deception]\label{p:deceive}
Deceiving a party about the terms of a transaction treats him merely as a means, unless he is pursuing a wrongful end that the deception is needed to frustrate.
\end{proposition}

\begin{proof}
Deceiving him uses his beliefs to bring about his participation on terms he has not accepted. He has a ground to object to being so used, namely that his agency is directed by a false picture of what he is agreeing to, and the ground does not depend on his pursuing a wrongful end unless he is. This is Proposition~3.2 of ETH-21-a applied to transactions; the exception is its Proposition~3.3.
\end{proof}

\begin{objection}
Proposition~\ref{p:deceive} is too strong. A seller who conceals a defect and then gives the buyer a larger discount than the defect warrants leaves him better off.
\end{objection}

\begin{reply}
The buyer is better off but has been made to accept terms other than those he believed he accepted. His ground for objecting concerns his authority over what he agrees to, not his welfare, and it is not outweighed by any ground the seller has, since the seller could have disclosed the defect and offered the same discount. The case shows that the formula of humanity, read through the word \enquote{merely}, protects agency rather than welfare \citep[4:429--430]{Kant1785}.
\end{reply}

\begin{objection}
Proposition~\ref{p:bargain} makes the permissibility of a bargain depend on the other party's alternatives, which the bargainer may not know.
\end{objection}

\begin{reply}
Permissibility depends on the facts; blameworthiness depends on what the agent could reasonably know. A bargainer ignorant of the other's desperation may act wrongly without being blameworthy, as in other cases of non-culpable ignorance. Ordinary commercial practice, in which parties typically have alternatives, makes such cases rare.
\end{reply}

\section{Conclusion}

The formula of humanity, read through the word \enquote{merely}, permits the use of persons that ordinary economic life requires and forbids its two characteristic corruptions: deception about terms and the exploitation of a lack of acceptable alternatives. Neither verdict requires that the parties care about one another. Both turn on whether the way a person is used is one he could reasonably reject.

\begin{thebibliography}{9}\small
\bibitem[Kant(1785)]{Kant1785} Kant, I. (1785). \emph{Groundwork of the Metaphysics of Morals}. In \emph{Practical Philosophy}, trans. M.~J. Gregor, Cambridge University Press, 1996.
\bibitem[O'Neill(1985)]{ONeill1985} O'Neill, O. (1985). Between consenting adults. \emph{Philosophy \& Public Affairs}, 14(3), 252--277.
\bibitem[Wertheimer(1996)]{Wertheimer1996} Wertheimer, A. (1996). \emph{Exploitation}. Princeton University Press.
\end{thebibliography}

\end{document}
